% Figure: cryogenic tank boil-off. Left axis, the radiative heat leak
% through a multilayer blanket falls as 1/(N+1) with the number of ideal
% shields; right, the same leak converted to a daily boil-off fraction of
% a stored cryogen. Plotted as heat leak (W) versus number of shields on
% log-linear axes, with the bare-gap and ten-layer points marked.
\begin{figure}[htbp]
\centering
\begin{tikzpicture}
\begin{axis}[
  width=12cm, height=7cm,
  ymode=log,
  xlabel={number of radiation shields $N$},
  ylabel={radiative heat leak $\dot{Q}$ (W)},
  xmin=0, xmax=30,
  ymin=0.2, ymax=30,
  axis lines=left,
  tick label style={font=\scriptsize},
  label style={font=\small},
  legend style={font=\scriptsize, at={(0.98,0.98)}, anchor=north east},
  legend cell align=left,
  legend entries={ideal stack $\dot{Q}_0/(N+1)$},
  clip=false,
]
% Bare-gap leak Q0 = 20 W (schematic). Ideal stack: Q0/(N+1).
\addplot[engblue, very thick, domain=0:30, samples=31] {20/(x+1)};
\addplot[only marks, engred, mark=*, mark size=2pt] coordinates {(0,20) (10,1.82)};
\node[font=\scriptsize, engred, anchor=south west] at (axis cs:0,20) {bare gap};
\node[font=\scriptsize, engred, anchor=south west] at (axis cs:10,1.82) {10 shields};
\end{axis}
\end{tikzpicture}
\caption[Heat leak through a multilayer cryogenic blanket]{An ideal stack
of $N$ floating radiation shields cuts the bare-gap heat leak by a factor
$1/(N+1)$, so ten shields drop the leak roughly elevenfold and the curve
flattens: each further shield returns less than the one before, and a real
blanket stops improving once the spacer netting that separates the shields
conducts across them faster than the added shield blocks. This is why a
practical cryogenic Dewar quotes a measured effective conductivity for the
whole blanket rather than an ideal shield count, and why the last decade of
heat leak is fought with better spacers and fewer penetrations rather than
with more layers. The leak fixes the boil-off: the escaping heat divided by
the latent heat of the cryogen is the mass lost per unit time, the number a
tank is bought and sized against.}
\label{fig:vol04ch05:cryostat-boiloff}
\end{figure}
